% !PRINTED 196
\weberpage{196}
%<<CONT
und wir betrachten unsere Aufgabe als gelöst, wenn von jeder
dieser Schaaren ein Repräsentant bestimmt ist.

Wir werden diese Freiheit später benutzen, um die gefundenen
Gruppen möglichst einfach darzustellen.

Die Determinante
\[
\alpha\delta - \beta\gamma = \Delta
\]
muss von Null verschieden sein, und wenn es die Einfachheit
verlangt, können wir sie immerhin $= 1$ annehmen, was wir bisweilen
thun werden.

Wir bezeichnen im Sinne der Gruppentheorie die Wiederholung
einer Substitution durch Exponenten: $\Theta, \Theta^2, \Theta^3, \ldots$,
worunter also nicht Potenzen, sondern immer wieder lineare
Substitutionen der Gruppe (1) zu verstehen sind. Die identische
Substitution $x = x$, die als die Einheit der Gruppe anzusehen
ist, wird auch mit $\Theta^0$ oder mit 1 bezeichnet.

Zwei inverse Substitutionen $\Theta, \Theta^{-1}$ können in der Form
dargestellt werden:
\[
\Theta = \begin{pmatrix} \alpha, & \beta \\ \gamma, & \delta \end{pmatrix}, \quad
\Theta^{-1} = \begin{pmatrix} \delta, & -\beta \\ -\gamma, & \alpha \end{pmatrix}.
\]

Da die Gruppe nach der Voraussetzung endlich ist, so hat
jedes ihrer Elemente einen bestimmten \gesperrt{Grad}, d. h. es giebt für
jedes $\Theta$ eine bestimmte kleinste positive Zahl $t$, für die $\Theta^t = 1$
ist. Diese Zahl $t$ muss ein Theiler von $n$ sein.

Für die Folge ist es von Wichtigkeit, für jede der Substitutionen
der Gruppe (ausgenommen die identische Substitution)
die Werthe der Variablen zu betrachten, die ihren Transformirten
gleich werden, also die Wurzeln der Gleichungen
\weq{5}{x = \Theta(x).}
Diese Werthe wollen wir die \gesperrt{Pole} der Substitution $\Theta$
nennen.

Die Gleichung (5) ist quadratisch und nimmt, wenn man
für $\Theta$ den Ausdruck (2) einsetzt, die Form an:
\weq{6}{\gamma x^2 + (\delta - \alpha) x - \beta = 0.}
Diese Gleichung hat zwei Wurzeln, die nur dann einander
gleich sind, wenn
\weq{7}{(\delta - \alpha)^2 + 4\beta\gamma = 0}
oder
\weq{8}{(\alpha + \delta)^2 = 4\Delta}
%CONT>>
